\section*{Propuesta Índice}

\setlist{nosep, leftmargin=*} 

\begin{enumerate}
  % --- CAPÍTULO 1 ---
  \item \textbf{Introducción}
  \begin{enumerate}[label*=\arabic*.]
    \item Historia de la energía eólica
    \item Contexto global y nacional de la energía eólica
    \item Desafíos actuales (incremento en el número de turbinas y tamaño de rotores=
    \item Objetivos generales y específicos
    \item Estructura de la tesis
  \end{enumerate}

  % --- CAPÍTULO 2 ---
  \item \textbf{Marco teórico}
  \begin{enumerate}[label*=\arabic*.] 
  	\item Dinámica de Fluidos Computacioanl (CFD)
  	\begin{enumerate}[label*=\arabic*.]
  		\item Ecuaciones de N-S
  		\item Métodos de Volúmenes Finitos
  		\item Modelado de Turbulencia: RANS y LES
  		\item OpenFOAM
  	\end{enumerate}
    	\item Capa Límite Atmosférica
    		\begin{enumerate}[label*=\arabic*.]
    		\item Características generales
    		\item Perfiles de velocidades de viento
    		\item Efecto de la temperatura
    		\item Atmospheric surface layer vs atmospheric boundary layer
    		\item CNBL
    		\end{enumerate}
    \item Aerodinámica de Aerogeneradores
    		\begin{enumerate}[label*=\arabic*.]
    			\item Teoría Unidimensional de Momento
    			\item Teoría BEM
    			\item Aerodinámica real de aerogeneradores
    			\item Efecto estela en Aerogeneradores
    			\item Interacción con la CNABL
    		\end{enumerate}
    	\item Modelos para la simulación de aergeneradores
    		\begin{enumerate}[label*=\arabic*.]
    		\item Modelo de Actuador Discal
    		\item Modelo de Actuador Lineal
    		\item Otros modelos
    		\end{enumerate}
    	\item Proper Orthogonal Decomposition
    	\item Phase Average
  \end{enumerate} 

  % --- CAPÍTULO 3 ---
  \item \textbf{Desarrollo de la biblioteca wtActuatorFoam}
  \begin{enumerate}[label*=\arabic*.]
  	\item Estado del arte en el modelo de actuador discal y lineal
    \item Modelos de Actuador Discal
    \begin{enumerate}[label*=\arabic*.]
    		\item uniform
    		\item airfoil
    		\item numeric
    		\item adaptive
    		\item analytic
		\item genalytic
    \end{enumerate}
    \item Modelos de Actuador Lineal
    \begin{enumerate}[label*=\arabic*.]
    		\item uniform
    		\item airfoil
    		\item numeric
    		\item adaptive
    		\item analytic
		\item genalytic
	\end{enumerate}
     \item{Features varios}
    	\begin{enumerate}[label*=\arabic*.]
    			\item Actuator mesh
    			\item Induction model
    			\item tipFactor
    			\item rootFactor
    			\item weighting parameters
    			\item save Level
    	\end{enumerate}
    	\item Análisis de Escalabilidad
    	\item Validaciones
    	\begin{enumerate}[label*=\arabic*.]
    		\item Actuadores discales - WESC
    		\item Actuadores lineales - TORQUE
    	\end{enumerate}
    	\item Caso de Aplicación: Análisis de rotación de la estela en casos de ABL - Paper José
	\item Conclusiones del capítulo
  \end{enumerate}

  % --- CAPÍTULO 4 ---
  \item \textbf{Influencia de las condiciones de entrada en la estela de aergeneradores flotantes}
  \begin{enumerate}[label*=\arabic*.]
    \item Introducción
    \begin{enumerate}[label*=\arabic*.]
      \item Contexto actual 
      \item Desafíos nuevos debido al movimiento de la plataforma 
      \item Efecto estela en turbinas flotantes
      \item Estado del arte
    \end{enumerate}

	\item Metodología
    \begin{enumerate}[label*=\arabic*.]
      \item Simulación de un disco poroso ante flujo uniforme laminar y turbulento
      \begin{enumerate}[label*=\arabic*.]
      	\item Mallado
      	\item Condiciones de borde
      	\item Modelo de LES
      	\item Esquema de simulaciones
      \end{enumerate}
      \item Simulación de un disco poroso ante flujo atmosférico a escala 
       \begin{enumerate}[label*=\arabic*.]
      	\item Mallado
      	\item Condiciones de borde
      	\item Modelo de LES
      	\item Esquema de simulaciones
      \end{enumerate}
      \end{enumerate}
  \item Análisis de Resultados
  	\begin{enumerate}[label*=\arabic*.]
  	  \item Análisis de flujos medios
	  \item Q-criterion
  	\item POD
 	\item Promediado en fase (phase average)
  	\end{enumerate}

 \item Conclusiones del capítulo
    \end{enumerate}

  % --- CAPÍTULO 5 ---
  \item \textbf{Simulaciones URANS con tendencias atmosféricas en OpenFOAM}
  \begin{enumerate}[label*=\arabic*.]
    \item Introducción
    \begin{enumerate}[label*=\arabic*.]
    	  \item ABL con incorporación de Temperatura vs limitación de la longitud de mezcla
      \item Incorporación de la temperatura en CFD - Hipótesis de Boussinesq
      \item Generación de condiciones reales
      \begin{enumerate}[label*=\arabic*.]
      	\item Generación de BC en precursor
      	\item Uso de tendencias acopladas desde WRF
      \end{enumerate}
      \item Estado del arte: SOWFA - CFDWind 2.4.0
      \item Modelo de turbulencia kEpsilon de Sogachev
    \end{enumerate}
	\item Migración del CFDWind-v2412
	\begin{enumerate}[label*=\arabic*.]
		\item Funciones de pared ya implementadas en OF 2412
		\item Funciones de pared reimplementadas en OF 2412
		\item Reimplementación del modelo de turbulencia
		\item Validación GABLS3
		\item Validación GABLS2
	\end{enumerate}	  
	\item Casos de Aplicación: RAWSON / Bahía Blanca
    \item Conclusiones del capítulo 
  \end{enumerate}

  % --- CAPÍTULO 6 ---
  \item \textbf{Síntesis y aportes finales}
  \begin{enumerate}[label*=\arabic*.]
    \item Avances sobre actuadores
    \item Avances sobre estelas de turbinas flotantes
    \item Avances sobre acoples WRF-CFD con URANS
    \item Líneas futuras de investigación
  \end{enumerate}

  \item \textbf{Conclusiones generales}

  \item \textbf{Referencias bibliográficas}
\end{enumerate}